\documentclass[11pt]{article}
\usepackage{preamble}
\renewcommand{\answer}[1]{}

\begin{document}

\ladderheading{AP Statistics \textperiodcentered\ Unit 1 \textperiodcentered\ Topic 1.10 \textperiodcentered\ B: 2026-09-25 \textperiodcentered\ E: 2026-10-02}{Who Chose Whom?}

\focusbanner{TODAY'S PROBLEM}

Suppose a school website posts the question, ``Should the cafeteria go cashless?'' Any visitor may answer. Of 412 responses, $71\%$ say Yes. That same week, a teacher randomly selects three homerooms and asks every student in those rooms. Of 68 responses, $44\%$ say Yes.\par\smallskip The website headline is ``Students back a cashless cafeteria.'' The teacher's headline is ``Students are split on going cashless.''\par
\begin{center}
\textbf{Two polls of student opinion}\par\smallskip
\begin{tabular}{|l|c|c|c|}
\hline
\textbf{Poll} & \textbf{$n$} & \textbf{Yes} & \textbf{Selection} \\ \hline
\textbf{Website} & 412 & $71\%$ & Visitors opt in \\ \hline
\textbf{Homerooms} & 68 & $44\%$ & 3 rooms random \\ \hline
\end{tabular}
\end{center}
\begin{firsttakebox}{FIRST TAKE --- before the video (5 min)}
Which headline would you trust more for what students at this school think? Name one detail behind your choice.
\workspace{0.9in}
\end{firsttakebox}

\begin{callout}{\IconBook\ TRUST STARTS WITH SELECTION (keep on the page)}{calloutgreen}
The \textbf{population} is the whole group we want to learn about; the \textbf{sample} is the group we collect data from. Ask how the data were collected, who was included, and who was left out. A representative sample reflects the population. Chance selection helps avoid favoring one part of the population. More responses do not by themselves make a sample representative.
\end{callout}

\newpage
\textbf{Finish after the video:}\par\smallskip

\dokbadge{1}~\textbf{(a)} Who is the population the headlines are trying to describe?\par
\workspace{0.7in}

\dokbadge{2}~\textbf{(b)} In two sentences, explain how students entered each poll and which students the website poll could miss.\par
\workspace{1.4in}

\dokbadge{3}~\textbf{(c)} Which percentage is more trustworthy for this school? Explain how selection supports your choice and why the larger response count does not settle it. Suggest one change that directly improves who enters the website sample.\par
\workspace{2.0in}

\begin{sentenceframebox}\raggedright\textbf{Frame:} I would trust $\blank[0.5in]\%$ more because the sample was formed by \blank[2.2in].\end{sentenceframebox}

\begin{sentenceframebox}\raggedright\textbf{Frame:} The 412 responses do not fix \blank[1.4in] because \blank[2.3in].\end{sentenceframebox}

\begin{summaryexitbox}{EXIT --- turn this in}
One thing the video changed about my first take: \hrulefill\par
\workspace{0.4in}
\end{summaryexitbox}

\end{document}
